\rtmtight
\begin{tblr}{width=\textwidth,colspec={|Q[c,m,2.1cm]|X[l,m]|},hlines,vlines,hspan=minimal,rowsep=1.5pt,colsep=3pt,row{1}={bg=headgray,font=\bfseries}}
\SetCell{c,m}\hfil\logo\hfil & \textbf{POLITEKNIK NEGERI MALANG}\par\textbf{JURUSAN TEKNOLOGI INFORMASI}\par\textbf{PROGRAM STUDI: D4 TEKNIK INFORMATIKA} \\
\SetCell[c=2]{l} \textbf{RENCANA TUGAS MAHASISWA (RTM) DAN RUBRIK PENILAIAN} & \\
Nama Mata Kuliah & Penjaminan Mutu Perangkat Lunak \\
Kode Mata Kuliah & RTI255006 \quad \textbf{sks / Jam perminggu:} 2 SKS / 4 Jam \quad \textbf{Semester:} 5 \\
Dosen Pengampu & \begin{enumerate}\item Dian Hanifudin Subhi, S.Kom., M.Kom.\item Dika Rizky Yunianto, S.Kom., M.Kom.\item M. Hasyim Ratsanjani, S.Kom., M.Kom.\end{enumerate} \\
\end{tblr}

\assessmentblock{Tugas 1: Analisis Standar ISO dan QA}{Case Method}{SCPMK0608-04201}{Mahasiswa menganalisis standar mutu perangkat lunak dan relevansinya pada kasus aplikasi.}{Menganalisis kasus, memilih standar/quality model, dan menyusun argumentasi tertulis.}{Laporan analisis PDF.}{Analisis masalah mutu perangkat lunak & 0.8 & Masalah mutu pada kasus diidentifikasi tepat dan dikaitkan dengan dampak terhadap pengguna, proses, dan produk. & Masalah mutu cukup tepat, tetapi dampak atau konteks kasus belum lengkap. & Masalah mutu tidak jelas atau tidak terkait kasus. \\ \\
Pemilihan standar dan quality model & 0.7 & ISO/IEC 25010 atau standar QA relevan dipilih, dijelaskan, dan dipakai untuk membaca kasus. & Standar yang dipilih cukup relevan, tetapi penerapannya pada kasus masih terbatas. & Tidak menggunakan standar mutu relevan atau salah menerapkannya. \\ \\
Argumentasi dan bukti laporan & 0.5 & Argumen runtut, berbasis evidence kasus, dan laporan mudah ditelusuri. & Argumen cukup runtut, tetapi evidence atau struktur laporan belum kuat. & Argumen lemah, tidak berbasis evidence, atau laporan sulit dipahami. \\}

\assessmentblock{Tugas 2: Identifikasi Tipe dan Level Pengujian}{Case Method}{SCPMK0608-04201}{Mahasiswa mengidentifikasi pendekatan, tipe, dan level pengujian pada studi kasus.}{Mengklasifikasikan kebutuhan uji, menentukan tipe/level uji, dan memberi alasan pemilihan.}{Laporan klasifikasi pengujian PDF.}{Klasifikasi pendekatan pengujian & 0.8 & Manual, automated, black-box, white-box, dan gray-box diklasifikasikan sesuai kebutuhan kasus. & Klasifikasi cukup tepat, tetapi beberapa pendekatan belum dibedakan jelas. & Klasifikasi keliru atau tidak sesuai kasus. \\ \\
Pemetaan tipe dan level pengujian & 0.7 & Unit, integration, system, acceptance, dan tipe uji lain dipetakan ke requirement dengan alasan logis. & Pemetaan cukup tepat, tetapi alasan atau cakupan requirement belum lengkap. & Tidak mampu memetakan level pengujian ke requirement. \\ \\
Rationale dan keterkaitan SDLC & 0.5 & Pilihan pengujian dijelaskan sesuai risiko, fase SDLC, dan tujuan verifikasi/validasi. & Rationale tersedia, tetapi hubungan dengan risiko atau SDLC masih umum. & Tidak ada rationale yang dapat dipertanggungjawabkan. \\}

\assessmentblock{Tugas 3: Test Plan, Test Scenario, dan Test Case}{Project Based Learning}{SCPMK0608-04202}{Mahasiswa menyusun dokumen perencanaan pengujian berdasarkan requirement.}{Menganalisis requirement, membuat test plan, scenario, test case, data uji, dan expected result.}{Dokumen test plan, scenario, dan test case.}{Kelengkapan test plan & 1.6 & Scope, objective, resource, schedule, risk, entry/exit criteria, dan strategy ditulis lengkap. & Test plan tersedia, tetapi beberapa elemen perencanaan belum detail. & Test plan tidak lengkap atau tidak dapat menjadi acuan pengujian. \\ \\
Kualitas scenario dan test case & 1.4 & Scenario, test step, test data, expected result, dan precondition jelas serta dapat dieksekusi. & Scenario/test case cukup dapat dieksekusi, tetapi detail data atau expected result kurang. & Scenario/test case tidak jelas atau tidak dapat dieksekusi. \\ \\
Traceability requirement & 1 & Setiap test case tertaut ke requirement dan cakupan pengujian dapat dilacak. & Traceability tersedia sebagian, tetapi belum mencakup semua requirement penting. & Tidak ada traceability yang memadai. \\}

\assessmentblock{Kuis 1: Konsep SQA dan Perencanaan Awal}{Kuis berbasis checkpoint PBL kelompok}{SCPMK0608-04201, SCPMK0608-04202}{Mahasiswa mengevaluasi kesiapan awal proyek kelompok melalui konteks QA, pemetaan standar mutu, ruang lingkup pengujian, risiko awal, dan traceability awal.}{Kelompok menyusun artefak checkpoint, mempresentasikan rationale QA, dan menjawab validasi dosen terhadap keputusan awal proyek.}{Dokumen checkpoint QA awal dan catatan validasi/presentasi kelompok.}{Konteks QA proyek & 4 & Konteks aplikasi, masalah mutu, stakeholder, dan quality attribute awal dijelaskan jelas. & Konteks QA cukup jelas, tetapi stakeholder atau quality attribute belum lengkap. & Konteks QA tidak jelas atau tidak terkait proyek. \\ \\
Pemetaan standar, scope, dan risiko & 3.5 & Standar mutu, testing scope, risiko awal, dan prioritas pengujian dipetakan logis. & Pemetaan tersedia, tetapi prioritas risiko atau scope masih kurang konsisten. & Pemetaan standar/scope/risiko tidak tepat. \\ \\
Traceability awal dan defense & 2.5 & Traceability awal dapat dilacak dan kelompok mampu menjawab validasi dosen dengan evidence. & Traceability tersedia sebagian dan defense cukup menjawab pertanyaan utama. & Traceability tidak jelas atau kelompok tidak mampu mempertahankan keputusan. \\}

\assessmentblock{Tugas 4: Review Dokumentasi dan Presentasi}{Peer Review}{SCPMK0608-04202}{Mahasiswa mengevaluasi dokumen pengujian dan mempresentasikan rekomendasi perbaikan.}{Melakukan peer review, mencatat temuan, merevisi dokumen, dan mempresentasikan hasil.}{Lembar review, dokumen revisi, dan bahan presentasi.}{Kedalaman peer review & 0.8 & Review menemukan masalah substansi pada test plan, scenario, test case, dan traceability. & Review menemukan beberapa masalah, tetapi masih banyak komentar umum. & Review dangkal atau hanya memeriksa format. \\ \\
Rekomendasi dan revisi artefak & 0.7 & Rekomendasi spesifik, dapat ditindaklanjuti, dan revisi artefak terlihat jelas. & Sebagian rekomendasi ditindaklanjuti, tetapi prioritas atau bukti revisi kurang. & Tidak ada rekomendasi bermakna atau revisi tidak dilakukan. \\ \\
Kualitas presentasi hasil review & 0.5 & Presentasi menjelaskan temuan, dasar penilaian, keputusan revisi, dan dampaknya pada kualitas pengujian. & Presentasi cukup jelas, tetapi dasar penilaian atau dampak revisi belum kuat. & Presentasi tidak menjelaskan hasil review secara memadai. \\}

\assessmentblock{Tugas 5: Monitoring, Control, Completion, dan Report}{Project Based Learning}{SCPMK0608-04202}{Mahasiswa menyusun monitoring, control, completion, dan laporan pengujian.}{Menentukan metrik, memonitor progres uji, mencatat risiko, dan menyusun test report.}{Dokumen test monitoring dan test report.}{Metrik monitoring dan control & 0.8 & Progress, defect, risk, coverage, dan action item dimonitor dengan metrik yang sesuai. & Metrik tersedia, tetapi belum lengkap atau belum dipakai untuk control. & Monitoring tidak berbasis metrik atau status pengujian tidak jelas. \\ \\
Test completion dan laporan & 0.7 & Completion criteria, summary eksekusi, defect summary, dan status release disusun jelas. & Laporan tersedia, tetapi completion criteria atau defect summary belum lengkap. & Laporan tidak menggambarkan hasil pengujian. \\ \\
Rekomendasi tindak lanjut & 0.5 & Rekomendasi berbasis data monitoring, risiko, dan defect yang ditemukan. & Rekomendasi cukup relevan, tetapi belum kuat berbasis data. & Rekomendasi tidak operasional atau tidak terkait hasil pengujian. \\}

\assessmentblock{Tugas 6: Functional Testing}{Hands-on Practice}{SCPMK0608-04203}{Mahasiswa menerapkan functional testing untuk memverifikasi requirement aplikasi.}{Menyusun skenario functional test, menjalankan pengujian, dan mendokumentasikan evidence.}{Dokumen functional test dan evidence hasil uji.}{Coverage requirement fungsional & 1.2 & Requirement utama tercakup oleh skenario fungsional positif dan negatif. & Coverage cukup, tetapi beberapa requirement atau negative case belum diuji. & Coverage tidak memadai atau tidak terkait requirement. \\ \\
Kualitas test data dan eksekusi & 1.1 & Data uji, step, expected result, actual result, dan status eksekusi lengkap. & Eksekusi tercatat, tetapi data uji atau actual result belum detail. & Eksekusi tidak terdokumentasi atau hasil tidak dapat diverifikasi. \\ \\
Evidence defect dan analisis & 0.7 & Defect dicatat dengan bukti, severity, langkah reproduksi, dan analisis dampak. & Defect/evidence tersedia, tetapi severity atau analisis dampak belum lengkap. & Tidak ada evidence defect atau analisis hasil uji. \\}

\assessmentblock{Tugas 7: User Acceptance Testing}{Role Playing}{SCPMK0608-04203}{Mahasiswa merancang dan menjalankan UAT berdasarkan user story dan acceptance criteria.}{Menyusun UAT plan, menjalankan simulasi user, dan mencatat hasil eksekusi.}{Dokumen UAT plan dan execution record.}{Kesesuaian UAT dengan user story & 1.2 & UAT plan, user story, acceptance criteria, dan peran user konsisten. & UAT cukup sesuai, tetapi acceptance criteria atau role user belum lengkap. & UAT tidak berangkat dari kebutuhan pengguna. \\ \\
Simulasi dan pencatatan hasil UAT & 1.1 & Sesi UAT dijalankan, feedback user dicatat, status diterima/ditolak jelas. & Simulasi dilakukan, tetapi catatan feedback atau status acceptance belum lengkap. & UAT tidak dieksekusi atau hasil tidak terdokumentasi. \\ \\
Rekomendasi perbaikan berbasis UAT & 0.7 & Temuan UAT diterjemahkan menjadi rekomendasi prioritas yang dapat ditindaklanjuti. & Rekomendasi tersedia, tetapi prioritas atau keterkaitan temuan belum kuat. & Tidak ada rekomendasi berbasis hasil UAT. \\}

\assessmentblock{UTS: Strategi dan Dokumentasi Pengujian}{UTS berbasis milestone review PBL kelompok}{SCPMK0608-04201, SCPMK0608-04202}{Mahasiswa mengevaluasi kemajuan tengah semester proyek kelompok melalui strategi QA, test plan, test scenario, test case, traceability matrix, dan monitoring/control plan.}{Kelompok menyerahkan dokumen milestone, mempresentasikan artefak, melakukan defense, dan merevisi berdasarkan umpan balik.}{Dokumen milestone UTS, traceability matrix, monitoring/control plan, dan catatan defense kelompok.}{Strategi QA dan test plan milestone & 6 & Strategi QA, scope, risk, schedule, resource, dan test plan milestone lengkap dan realistis. & Strategi tersedia, tetapi beberapa aspek scope, risk, atau resource belum rinci. & Strategi QA tidak realistis atau tidak sesuai proyek. \\ \\
Konsistensi scenario, test case, dan traceability & 5.3 & Scenario, test case, data uji, expected result, dan traceability matrix saling konsisten. & Dokumen cukup konsisten, tetapi beberapa requirement belum terlacak. & Dokumen tidak konsisten atau tidak dapat dipakai untuk eksekusi. \\ \\
Monitoring plan dan defense kelompok & 3.7 & Monitoring/control plan jelas dan kelompok mampu mempertahankan keputusan pengujian dengan evidence. & Defense cukup baik, tetapi monitoring plan atau evidence masih kurang kuat. & Kelompok tidak mampu mempertahankan dokumen milestone. \\}

\assessmentblock{Tugas 8: E2E Testing}{Hands-on Practice}{SCPMK0608-04203}{Mahasiswa mengimplementasikan end-to-end testing pada alur utama aplikasi.}{Menentukan alur kritis, membuat script E2E, menjalankan test, dan melampirkan evidence.}{Script E2E dan laporan hasil eksekusi.}{Pemilihan alur E2E kritis & 1.2 & Alur bisnis utama, precondition, data, dan expected outcome dipilih sesuai risiko aplikasi. & Alur E2E cukup relevan, tetapi data atau expected outcome belum lengkap. & Alur E2E tidak kritis atau tidak merepresentasikan penggunaan nyata. \\ \\
Kualitas script dan assertion & 1.1 & Script E2E stabil, readable, memakai selector tepat, assertion kuat, dan menangani setup/teardown. & Script berjalan, tetapi selector, assertion, atau setup masih rapuh. & Script tidak berjalan atau assertion tidak memvalidasi hasil. \\ \\
Evidence eksekusi dan analisis & 0.7 & Hasil run, screenshot/log, failure analysis, dan rekomendasi perbaikan terdokumentasi. & Evidence run tersedia, tetapi analisis failure masih terbatas. & Tidak ada evidence eksekusi atau hasil tidak dianalisis. \\}

\assessmentblock{Tugas 9: API Testing Suite}{Hands-on Practice}{SCPMK0608-04203}{Mahasiswa membuat API testing suite manual dan otomatis.}{Membuat request, environment, assertion, negative case, dan automated run.}{Postman/Newman collection dan laporan eksekusi.}{Desain request dan environment & 1.2 & Collection, endpoint, method, parameter, body, header, token, dan environment disusun benar. & Request berjalan, tetapi environment, token, atau parameter belum konsisten. & Request tidak sesuai kontrak API atau sulit dijalankan ulang. \\ \\
Assertion dan negative case & 1.1 & Assertion status, schema, payload, business rule, serta negative case mencakup kontrak API. & Assertion tersedia, tetapi negative case atau validasi schema belum lengkap. & Assertion tidak memadai atau hanya mengecek response dasar. \\ \\
Automation dan laporan eksekusi & 0.7 & Newman/automation berjalan, laporan hasil tersedia, dan failure dianalisis. & Automation berjalan sebagian, tetapi laporan atau analisis failure belum lengkap. & Tidak ada automation/evidence eksekusi yang valid. \\}

\assessmentblock{Tugas 10: Unit dan Integration Testing}{Coding Practice}{SCPMK0608-04203}{Mahasiswa membuat unit test dan integration test pada fungsi aplikasi.}{Menulis test code, menjalankan test suite, dan mengevaluasi hasil test.}{Kode test dan laporan hasil eksekusi.}{Desain unit test fungsi kritis & 1.2 & Unit test mencakup fungsi kritis, boundary case, error case, dan isolasi dependency. & Unit test berjalan, tetapi boundary/error case masih terbatas. & Unit test tidak relevan atau tidak mencakup fungsi kritis. \\ \\
Integration test antar komponen & 1.1 & Integration test memverifikasi alur data antar modul/service/database sesuai skenario. & Integration test tersedia, tetapi cakupan alur atau setup data belum lengkap. & Integration test tidak berjalan atau tidak menguji integrasi nyata. \\ \\
Kualitas test code dan TDD & 0.7 & Test code rapi, repeatable, readable, dan menunjukkan siklus TDD/refactor. & Test code cukup rapi, tetapi repeatability atau bukti TDD masih terbatas. & Test code sulit dijalankan ulang atau tidak terstruktur. \\}

\assessmentblock{Tugas 11: Review Test Code dan Coverage}{Peer Assessment}{SCPMK0608-04203}{Mahasiswa mengevaluasi kualitas test code dan test coverage.}{Melakukan code review, membaca coverage report, dan menyusun rekomendasi refactoring.}{Laporan review test code dan coverage.}{Ketepatan code review test & 1.2 & Review mengidentifikasi masalah readability, isolation, flaky test, duplication, dan assertion lemah. & Review menemukan beberapa masalah, tetapi belum menyentuh risiko test penting. & Review dangkal atau tidak berbasis kode test. \\ \\
Analisis coverage & 1.1 & Coverage report dibaca tepat, gap fungsi kritis diidentifikasi, dan tidak hanya mengejar angka coverage. & Coverage dianalisis, tetapi gap fungsi kritis belum lengkap. & Coverage tidak dianalisis atau disimpulkan keliru. \\ \\
Rekomendasi refactoring test & 0.7 & Rekomendasi refactoring spesifik, prioritas jelas, dan meningkatkan maintainability test suite. & Rekomendasi cukup relevan, tetapi masih umum atau tidak diprioritaskan. & Rekomendasi tidak dapat ditindaklanjuti. \\}

\assessmentblock{Tugas 12: Analisis Non-Functional Testing}{Case Method}{SCPMK0608-04204}{Mahasiswa menganalisis kebutuhan non-functional testing berdasarkan quality attributes.}{Mengidentifikasi quality attributes, prioritas risiko, dan strategi non-functional testing.}{Laporan analisis non-functional testing.}{Identifikasi quality attributes & 1.2 & Performance, security, usability, reliability, maintainability, atau atribut relevan dipilih sesuai konteks. & Atribut kualitas cukup relevan, tetapi beberapa prioritas belum didukung konteks. & Atribut kualitas tidak relevan atau tidak teridentifikasi. \\ \\
Strategi non-functional testing & 1.1 & Strategi uji, metrik, tools, threshold, dan risiko dirancang sesuai kebutuhan pengguna. & Strategi tersedia, tetapi metrik, tools, atau threshold belum lengkap. & Strategi tidak operasional atau tidak menguji atribut kualitas. \\ \\
Prioritas risiko dan rekomendasi & 0.7 & Risiko non-fungsional diprioritaskan dan rekomendasi pengujian jelas. & Prioritas risiko cukup jelas, tetapi rekomendasi masih umum. & Tidak ada prioritas risiko atau rekomendasi tidak relevan. \\}

\assessmentblock{Tugas 13: Performance Testing Scenario}{Hands-on Practice}{SCPMK0608-04204}{Mahasiswa membuat dan menjalankan skenario performance testing.}{Mendesain skenario beban, menjalankan tools, membaca metrik, dan membuat rekomendasi.}{Skenario performance testing dan laporan metrik.}{Desain skenario beban & 1.2 & Jumlah user, ramp-up, workload, data, SLA, dan target metrik dirancang realistis. & Skenario beban tersedia, tetapi SLA, data, atau workload belum lengkap. & Skenario tidak realistis atau tidak dapat dijalankan. \\ \\
Eksekusi tools dan pengumpulan metrik & 1.1 & JMeter/k6 dijalankan benar, metrik response time, throughput, error rate, dan resource dicatat. & Tools berjalan, tetapi beberapa metrik atau konfigurasi belum lengkap. & Tools tidak berjalan atau metrik tidak valid. \\ \\
Analisis bottleneck dan rekomendasi & 0.7 & Bottleneck dianalisis dari metrik dan rekomendasi teknis disusun jelas. & Analisis tersedia, tetapi penyebab bottleneck atau rekomendasi masih umum. & Hasil performa tidak dianalisis. \\}

\assessmentblock{PBL/Proyek Akhir: Strategi QA Komprehensif}{Project Based Learning}{SCPMK0608-04201, SCPMK0608-04202, SCPMK0608-04203, SCPMK0608-04204}{Mahasiswa mengintegrasikan strategi QA lengkap pada aplikasi studi kasus.}{Menganalisis kebutuhan mutu, menyusun strategi QA, menjalankan functional dan non-functional testing, serta menyusun laporan akhir.}{Laporan proyek, artefak pengujian, evidence, repository/script, dan presentasi.}{Integrasi strategi QA proyek & 10 & Quality model, risk, test strategy, planning, functional, dan non-functional testing terintegrasi pada proyek. & Strategi cukup terintegrasi, tetapi sebagian risiko atau jenis uji belum lengkap. & Strategi QA terfragmentasi atau tidak sesuai proyek. \\ \\
Implementasi dan evidence pengujian & 8.8 & Functional, API/E2E/unit/integration, UAT, dan non-functional testing dijalankan dengan evidence lengkap. & Pengujian berjalan sebagian besar, tetapi beberapa evidence atau cakupan masih kurang. & Pengujian utama tidak berjalan atau evidence tidak memadai. \\ \\
Laporan, rekomendasi, dan presentasi & 6.2 & Laporan profesional, defect/metric dianalisis, rekomendasi actionable, dan presentasi meyakinkan. & Laporan dan presentasi cukup jelas, tetapi rekomendasi atau analisis belum mendalam. & Laporan tidak menyimpulkan mutu aplikasi atau rekomendasi tidak berguna. \\}

\assessmentblock{UAS: Evaluasi Mutu dan Rekomendasi}{UAS berbasis validasi akhir PBL kelompok}{SCPMK0608-04203, SCPMK0608-04204}{Mahasiswa memvalidasi hasil akhir proyek kelompok melalui demonstrasi implementasi QA, evidence pengujian fungsional dan non-fungsional, analisis defect, serta rekomendasi peningkatan mutu.}{Kelompok mendemonstrasikan hasil pengujian, menjelaskan evidence, menganalisis defect, menyampaikan rekomendasi, dan melakukan presentasi akhir.}{Laporan akhir QA, evidence pengujian, defect log, rekomendasi mutu, repository/script, dan bahan presentasi akhir.}{Validasi evidence akhir & 6 & Evidence functional dan non-functional testing lengkap, dapat ditelusuri, dan sesuai scope proyek. & Evidence tersedia, tetapi traceability atau cakupan scope belum lengkap. & Evidence tidak cukup untuk menilai mutu proyek. \\ \\
Analisis defect dan mutu perangkat lunak & 5.3 & Defect, severity, root cause, quality attribute, dan dampak pengguna dianalisis tepat. & Analisis defect cukup tepat, tetapi root cause atau dampak mutu belum lengkap. & Analisis defect lemah atau tidak berbasis data. \\ \\
Rekomendasi dan defense akhir & 3.7 & Rekomendasi peningkatan mutu prioritas, realistis, dan dipertahankan kuat saat presentasi akhir. & Rekomendasi cukup relevan, tetapi prioritas atau defense belum kuat. & Rekomendasi tidak operasional atau tidak mampu dipertahankan. \\}

\begin{tblr}{width=\textwidth,colspec={|Q[l,m,3.2cm]|X[l,m]|},hlines,vlines,hspan=minimal,row{1-3}={bg=headgray},rowsep=1.5pt,colsep=3pt}
Jadwal Pelaksanaan: & Minggu 1--17 sesuai RPS. Setiap evaluasi memiliki RTM dan rubrik multi-kriteria. \\
Lain-Lain Yang Diperlukan: & LMS, repositori/prototype/proyek, perangkat lunak sesuai mata kuliah, dan perangkat presentasi. \\
Pustaka: & Mengacu pustaka utama dan pendukung pada bagian RPS. \\
\end{tblr}
